\documentclass[aspectratio=169, 10pt, t]{beamer}
\usepackage[T5,T1]{fontenc}
\usepackage[utf8]{vietnam}
\usepackage{helvet}
\renewcommand{\familydefault}{\sfdefault}
\usepackage{beamerthemeMLIOT}   % custom theme (in same folder)
\setbeamerfont{title}{series=\bfseries} % Ép Beamer phải in đậm tiêu đề chính

% ---- Additional packages ----
\usepackage{amsmath, amssymb}
\usepackage{booktabs}
\usepackage{multicol}
\usepackage{caption}
\usepackage{hyperref}
\usepackage{array}
\usepackage{pifont}
\usepackage{ragged2e}
\usepackage{tikz}
\usetikzlibrary{shapes.geometric, arrows.meta, positioning}

% ---- Hyperlink colours ----
\hypersetup{
  colorlinks = true,
  urlcolor   = MLIOTaccent,
  linkcolor  = MLIOTblue,
  citecolor  = MLIOTblue
}

% ============================================================
%  PAPER METADATA – fill these in
% ============================================================
\title[NotebookForge]{%
  \textbf{\Large NotebookForge}\\[0.2em]
  \large Hệ thống Tự động Sinh và Kiểm định Chất lượng Notebook Thực hành Học máy bằng Hệ thống Đa tác nhân sử dụng CrewAI
}

\author{%
  \textbf{\textcolor{MLIOTaccent}{Nguyễn Xuân Huy}}%
  \texorpdfstring{\textcolor{MLIOTgold}{\textsuperscript{*a}}}{*a}\textcolor{MLIOTnavy}{,~}\\%
  \textbf{\textcolor{MLIOTnavy}{Nguyễn Xuân Hoàng}}%
  \texorpdfstring{\textcolor{MLIOTgold}{\textsuperscript{a}}}{a}\textcolor{MLIOTnavy}{,~}%
  \textbf{\textcolor{MLIOTnavy}{Nguyễn Hoàng Nam}}%
  \texorpdfstring{\textcolor{MLIOTgold}{\textsuperscript{a}}}{a}\textcolor{MLIOTnavy}{,~}\\% 
  \textbf{\textcolor{MLIOTnavy}{Bùi Phong Hợp}}%
  \texorpdfstring{\textcolor{MLIOTgold}{\textsuperscript{a}}}{a}\textcolor{MLIOTnavy}{,~}%
  \textbf{\textcolor{MLIOTnavy}{Nguyễn Minh Trí}}%
  \texorpdfstring{\textcolor{MLIOTgold}{\textsuperscript{a}}}{a}\textcolor{MLIOTnavy}{~}%
}

\setAffiliations{%
  \textsuperscript{a} Nhóm 13: UTÚ
}

\setCorrespondingAuthor{%
  *Email nhóm trưởng: \href{mailto:author@mail.com}{author@mail.com}%
}

\setPaperID{MLIOT-2026-XXX}

\institute{%
  Machine Learning \& IoT Lab,
  Ho Chi Minh City University of Technology (HCMUT)%
}

\date{August 2026}

% ============================================================
\begin{document}
% ============================================================

% ------------------------------------------------------------
%  SLIDE 1 – Title Slide
% ------------------------------------------------------------
\begin{frame}[plain, c]
  \titlepage
\end{frame}

% ------------------------------------------------------------
%  SLIDE 2 – Outline
% ------------------------------------------------------------
\begin{frame}{Nội dung chính}
  \begin{columns}[T]
    \begin{column}{0.55\textwidth}
      \setbeamerfont{enumerate item}{size=\large}
      
      \begin{enumerate}
        \setlength{\itemsep}{0.8em}
        \item {\large Đặt vấn đề và giải pháp} 
        \item {\large Multi-agent Pipeline} 
        \item {\large Demo và Evaluation} 
        \item {\large Conclusion và Mentor Q\&A} 
        \item {\large Conclusions} 
        \item {\large References} 
      \end{enumerate}
    \end{column}
    \begin{column}{0.42\textwidth}
      \begin{block}{Contact Information}
        \footnotesize
        \textbf{\textcolor{MLIOTnavy}{Address:}}~403.1 BK.B6, HCMUT Campus 2\\[0.35em]
        \textbf{\textcolor{MLIOTnavy}{Email:}}\\
        mlandiotlab@gmail.com\\[0.35em]
        \textbf{\textcolor{MLIOTnavy}{Website:}}~mliotlab.com\\[0.35em]
        \textbf{\textcolor{MLIOTnavy}{Facebook:}}~facebook.com/hcmut.ml.iot.lab\\[0.35em]
        \textbf{\textcolor{MLIOTnavy}{YouTube:}}~youtube.com/@mliotlab
      \end{block}
    \end{column}
  \end{columns}
\end{frame}

% ------------------------------------------------------------
%  SLIDE 3 – Introduction
% ------------------------------------------------------------
\begin{frame}{Đặt vấn đề và giải pháp}
  \begin{columns}[T]
    \begin{column}{0.52\textwidth}
      \begin{block}{Tổng quan dự án \& bối cảnh thực tế}
        \begin{itemize}
          \item \textbf{Tên đề tài:} NotebookForge - 
          Hệ thống Tự động Sinh và Kiểm định Chất lượng 
          Notebook Thực hành Học máy bằng Hệ thống Đa tác nhân.
          \item \textbf{Thách thức của người học AI/ML:} Người mới bắt 
          đầu gặp khó khăn khi tìm kiếm bài tập thực hành 
          cá nhân hóa theo đúng trình độ, thời lượng rảnh 
          và mục tiêu cá nhân; tự học dễ nản vì thiếu lời giải chi tiết.
          \item \textbf{Hạn chế của LLM hiện tại:} Notebook 
          do ChatGPT/Claude thông thường sinh ra hay bị lỗi 
          không chạy được (runtime error), dễ bị rò rỉ dữ liệu 
          (data leakage), thiếu tính cấu trúc sư phạm, hoặc tự 
          bị đắt thông tin lý thuyết không chuẩn xác (hallucination).
        \end{itemize}
      \end{block}
    \end{column}
    \begin{column}{0.45\textwidth}
      \begin{block}{Literature Review Highlights}
        \begin{itemize}
          \item Author et al.\ (Year): Key finding 1.\cite{ref1}
          \item Author et al.\ (Year): Key finding 2.\cite{ref2}
        \end{itemize}
        \vspace{0.5em}
        \textbf{\textcolor{MLIOTcoral}{Research Gap:}}
        \textit{Describe gap identified.}
      \end{block}
    \end{column}
  \end{columns}
\end{frame}

% ------------------------------------------------------------
%  SLIDE 4 – Giải pháp NotebookForge với 3 mục tiêu chính
% ------------------------------------------------------------
\begin{frame}{Đặt vấn đề và giải pháp}
  Giải pháp \textbf{NotebookForge} được đặt ra với 3 mục tiêu chính:
  \begin{methodcols}
    \methodcol{\ding{202}~Cá nhân hóa lộ trình (Personalization)}{%
      \begin{itemize}
        \item Không sinh các bài học chung chung. 
        \item Hệ thống đánh giá năng lực thực tế qua bộ Quiz 5 câu 
        để xác định chính xác trình độ (\texttt{level\_final}) và dung lượng bài dựa theo thời lượng học do người dùng đăng ký (\texttt{duration\_minutes}).
      \end{itemize}%
    }
    \methodcol{\ding{203}~Tự thực thi \& sửa lỗi (Sandbox \& Self-correction)}{%
      \begin{itemize}
        \item Code không chỉ nằm trên giấy. Tất cả notebook được đưa vào môi trường cách ly (Sandbox) để chạy thử (Run-all). 
        \item Nếu phát hiện lỗi Python, agent sẽ đọc vết lỗi (Traceback log) và tự sửa code (Retry Loop) đến khi chạy thành công 100\%.
      \end{itemize}%
    }
    \methodcol{\ding{204}~Tri thức neo giữ \& Kiểm định kép (Grounding \& Dual-Verification)}{%
      \begin{itemize}
        \item Nội dung lý thuyết được neo giữ vào Knowledge Base chuẩn (Scikit-Learn User Guide) để chống bịa đặt.
        \item Đầu ra bắt buộc phải vượt qua bộ kiểm định kép: 8 Hard Rules (luật cứng kỹ thuật) và LLM Judge (chấm điểm chất lượng sư phạm).
      \end{itemize}%
    }
  \end{methodcols}
\end{frame}

% ------------------------------------------------------------
%  SLIDE 5 – Multi-agent
% ------------------------------------------------------------
\begin{frame}{Multi-agent Pipeline}
  Sơ đồ pipeline tổng quan:
  \begin{flushleft}
      \begin{tikzpicture}[
        node distance = 0.4cm and 0.45cm,
        box/.style = {
          rectangle, 
          rounded corners=3pt, 
          draw=structure.fg, 
          fill=structure.fg!10, 
          thick, 
          minimum width=2.4cm, 
          minimum height=1cm, 
          align=center,
          text width=2.2cm, 
          font=\fontsize{8pt}{10pt}\selectfont\bfseries,
        },
        arrow/.style = {
          -{Stealth[scale=0.8]}, 
          thick, 
          draw=structure.fg
        }
      ]
        \node (obj1) [box] {Learner Profile};
        \node (obj2) [box, right=of obj1] {Research Agent};
        \node (obj3) [box, right=of obj2] {Curriculum Agent};
        \node (obj4) [box, right=of obj3] {Notebook Generator};
        \node (obj5) [box, right=of obj4] {Executor \& Verifier};

        \draw [arrow] (obj1) -- (obj2);
        \draw [arrow] (obj2) -- (obj3);
        \draw [arrow] (obj3) -- (obj4);
        \draw [arrow] (obj4) -- (obj5);
      \end{tikzpicture}
  \end{flushleft}
  \begin{itemize}
    \item Mỗi Agent đóng vai trò như một công nhân chuyên môn hóa trong dây chuyền sản xuất notebook tự động với các nhiệm vụ riêng biệt.
  \end{itemize}
\end{frame}

% ------------------------------------------------------------
%  SLIDE 6 – Frontend Layer
% ------------------------------------------------------------
\begin{frame}{Multi-agent Pipeline}
  Thu thập yêu cầu \& đánh giá trình độ (Frontend Layer)
  \begin{methodcols}
    \methodcol{\ding{202}~Form cấu hình}{%
      Người dùng chọn:
        \begin{itemize}
          \item Topic (Logistic Regression, Decision Tree, K-Means Clustering)
          \item Level tự đánh giá (Beginner / Intermediate)
          \item Thời lượng bài học (\texttt{duration\_minutes}: 60--120 phút)
          \item Số lượng bài tập (1--5)
        \end{itemize}
      để xác định \texttt{level\_final} và \texttt{duration\_minutes}.
    }
    \methodcol{\ding{203}~Quiz 5 câu cho từng topic}{%
      \begin{itemize}
        \item Căn chỉnh trình độ thực tế. 
        \item Nếu người dùng chọn "Intermediate" nhưng làm sai nhiều câu Quiz, hệ thống tự động hạ chỉ số \texttt{level\_final} về đúng năng lực thực tế để tránh bài học quá sức.
      \end{itemize}%
    }
    \methodcol{\ding{204}~Cơ chế UI Fallback}{%
      Giao diện tự phục hồi và hiển thị thông báo trạng thái mượt mà nếu Server backend xử lý chậm hoặc gặp sự cố mạng.
    }
  \end{methodcols}
\end{frame}

% ------------------------------------------------------------
%  SLIDE 7 – Research Agent
% ------------------------------------------------------------
\begin{frame}{Multi-agent Pipeline}
  Research Agent: Cơ sở kiến thức \& Truy xuất dữ liệu
  \begin{itemize}
    \item \textbf{Knowledge Base (.md):} Xây dựng bộ tri thức chuẩn hóa từ tài liệu chính thống (Scikit-Learn User Guide), đóng vai trò "sách giáo khoa" cố định cho AI tham chiếu.
    \item \textbf{Luồng truy xuất 2 giai đoạn (Two-Stage Retrieval):}
      \begin{itemize}
        \item \textbf{Giai đoạn 1 (BM25 Retrieval):} Thuật toán tìm kiếm theo tần suất từ khóa chính xác để rút trích các đoạn văn bản lý thuyết liên quan từ Knowledge Base.
        \item \textbf{Giai đoạn 2 (LLM Re-ranking):} LLM sắp xếp và chọn lọc lại các khái niệm dựa trên thời lượng \texttt{duration\_minutes} (nhập từ UI).
        \begin{itemize}
          \item Bài 30 phút: Giữ lại 2--3 khái niệm cốt lõi.
          \item Bài 60 phút: Mở rộng thêm khái niệm chuyên sâu và bài tập nâng cao.
        \end{itemize}
      \end{itemize}
    \item \textbf{Fallback Web Search (Wikipedia API):} Khi Knowledge Base nội bộ không phủ hết khái niệm, hệ thống tự gọi Wikipedia API và áp dụng bộ lọc token hit ratio $\ge 0.5$ (chỉ giữ bài viết có tỷ lệ trùng khớp từ khóa từ 50\% trở lên để loại bỏ thông tin rác).
  \end{itemize}
\end{frame}

% ------------------------------------------------------------
%  SLIDE 8 – Curriculum & Generator Agent
% ------------------------------------------------------------
\begin{frame}{Multi-agent Pipeline}
  Curriculum \& Generator Agent: Sinh lộ trình \& Đóng gói Notebook
  \begin{itemize}
    \item \textbf{Curriculum Agent:} Phân tách luồng xử lý phù hợp cho Học có giám sát (Supervised) hoặc Không giám sát (Unsupervised).
    \item \textbf{Cấu trúc Notebook chuẩn 4 khối:}
      \begin{enumerate}
        \item \textbf{Tổng quan bài học:} Mục tiêu, lý thuyết nền tảng rút trích từ ResearchBundle.
        \item \textbf{Chuẩn bị dữ liệu \& EDA:} Nạp các bộ dữ liệu thực tế (Heart Failure, Red Wine Quality, Mall Customers...) và chèn sẵn các đoạn code trực quan hóa.
        \item \textbf{Nội dung \& Bài tập thực hành:} Phân chia tập Train/Test đúng chuẩn sư phạm, tuyệt đối không tiền xử lý dữ liệu trước khi split để tránh rò rỉ dữ liệu (Data Leakage).
        \item \textbf{Kiểm tra tự động:} Chèn sẵn các đoạn code dạng \texttt{assert} để người học tự kiểm tra đáp án.
      \end{enumerate}
  \end{itemize}
\end{frame}

% ------------------------------------------------------------
%  SLIDE 9 – Executor, Verifier & Retry Loop
% ------------------------------------------------------------
\begin{frame}{Multi-agent Pipeline}
  Executor, Verifier \& Vòng lặp Tự sửa lỗi
  \begin{itemize}
    \item \textbf{Executor Sandbox:} Môi trường chạy code ảo cách ly hoàn toàn, nhận file \texttt{.ipynb} và thực thi từ trên xuống dưới để bắt chính xác các lỗi runtime.
    \item \textbf{Bộ kiểm định kép (Verifier):}
      \begin{itemize}
        \item \textbf{8 Hard Rules (Quy tắc cứng):} Kiểm tra lỗi kỹ thuật bằng code (0 Execution Error, đúng định dạng JSON Schema, không Data Leakage, chứa đủ cell Todo/Assert...).
        \item \textbf{LLM Judge (Chấm điểm mềm):} Đánh giá chất lượng bài học trên 4 tiêu chí sư phạm (Khả năng thực thi, Độ chuẩn xác lý thuyết, Độ phù hợp trình độ, Tính logic sư phạm).
      \end{itemize}
    \item \textbf{Cơ chế Retry Loop:} Khi Verifier phát hiện lỗi, hệ thống không bỏ cuộc mà trích xuất vết lỗi dạng \texttt{[CELL n]... FIX:...} đưa lại cho Generator để sửa đúng ô code bị lỗi (tối đa 3 lần thử).
  \end{itemize}
\end{frame}

% ------------------------------------------------------------
%  SLIDE 10 – Demo
% ------------------------------------------------------------
\begin{frame}{Demo \& Evaluation}
  \begin{block}{Minh họa vận hành hệ thống}
    \begin{itemize}
      \item Quá trình tương tác trên giao diện Frontend: Chọn topic, thực hiện bài Quiz căn chỉnh trình độ.
      \item Minh họa quá trình sinh lộ trình tự động qua các Agent.
      \item Trực quan hóa vòng lặp tự phát hiện lỗi runtime và tự động sửa code (Retry Loop) trong Executor Sandbox.
      \item Sản phẩm đầu ra hoàn chỉnh: Notebook Jupyter đạt 100\% khả năng thực thi.
    \end{itemize}
  \end{block}
\end{frame}

% ------------------------------------------------------------
%  SLIDE 11 – Evaluation & Benchmark Metrics
% ------------------------------------------------------------
\begin{frame}{Demo \& Evaluation}
  Kết quả Kỹ thuật \& Benchmark Metrics
  \begin{itemize}
    \item \textbf{Golden Set (Bộ dữ liệu chuẩn):} Tập hợp 20 trường hợp kiểm thử thiết kế sẵn đại diện cho 3 chủ đề và 2 cấp độ trình độ khác nhau để đánh giá toàn diện hệ thống.
    \item \textbf{Các chỉ số đo lường chính:}
      \begin{itemize}
        \item \textbf{Execution Pass Rate:} Tỷ lệ Notebook mà 100\% các cell code chạy mượt mà từ trên xuống dưới, không phát sinh lỗi Python nào (0 Runtime Error).
        \item \textbf{Final PASS Rate:} Tỷ lệ Notebook vượt qua cả 8 Hard Rules kỹ thuật VÀ đạt điểm chất lượng sư phạm từ LLM Judge $\ge 4/5$.
        \item \textbf{Average LLM Rubric Score:} Điểm số chất lượng trung bình (thang điểm 5) do LLM Judge chấm dựa trên 4 tiêu chí sư phạm.
        \item \textbf{Cost \& Time per Case:} Chi phí API trung bình cho mỗi bài học ($\le 0.30\$$/case) và thời gian trung bình để hoàn tất vòng sinh \& kiểm thử.
      \end{itemize}
  \end{itemize}
\end{frame}

% ------------------------------------------------------------
%  SLIDE 12 – Conclusion & Mentor Q&A
% ------------------------------------------------------------
\begin{frame}{Conclusion \& Mentor Q\&A}
  \begin{columns}[T]
    \begin{column}{0.48\textwidth}
      \begin{block}{Tóm tắt đóng góp}
        \begin{itemize}
          \item Đề xuất kiến trúc Multi-agent tự động sinh Notebook thực hành ML cá nhân hóa.
          \item Tích hợp cơ chế kiểm định kép (Hard Rules + LLM Judge) đảm bảo độ chính xác kỹ thuật và giá trị sư phạm.
          \item Xây dựng vòng lặp tự sửa lỗi (Retry Loop) trong môi trường Sandbox giúp duy trì tỷ lệ thực thi code thành công cao.
        \end{itemize}
      \end{block}
    \end{column}
    
    \begin{column}{0.48\textwidth}
      \begin{block}{Kế hoạch hoàn thiện \& Q\&A}
        \begin{itemize}
          \item Mở rộng Knowledge Base sang các chủ đề Deep Learning và Computer Vision.
          \item Tối ưu hóa thời gian sinh bài và chi phí gọi LLM API.
          \item Thử nghiệm thực tế với nhóm người dùng sinh viên để thu thập feedback.
        \end{itemize}
      \end{block}
    \end{column}
  \end{columns}
\end{frame}

% ------------------------------------------------------------
%  SLIDE 13 – Thank You / Acknowledgements
% ------------------------------------------------------------
\thankyouslide{%
  \textit{The authors gratefully acknowledge Machine Learning \& IoT Lab, HCMUT for their support and guidance.}
}

% ============================================================
\end{document}
% ============================================================